\section{Introducción: qué es la visión por computadora}

La visión es uno de los sentidos más importantes del ser humano. Nos permite detectar e interpretar emociones y expresiones faciales, desplazarnos por el mundo, manipular objetos e interactuar con el entorno. Alexander Amini propuso al comienzo de esta clase una definición concisa: la visión por computadora (Computer Vision) consiste en dotar a las computadoras de la capacidad de "saber qué hay dónde con solo mirar" (\textit{to know what is where by looking}).

\begin{figure}[H]
\centering
\includegraphics[width=0.85\textwidth]{slides-images/slide-02.jpg}
\caption{Definición central de la visión por computadora: "To know what is where by looking." \protect\footnotemark}
\end{figure}
\footnotetext{Diapositiva 2.}

La visión no consiste solamente en reconocer objetos de manera estática, sino también en comprender los cambios dinámicos de una escena. Por ejemplo, cuando vemos la fotografía de una calle de una ciudad, no solo identificamos autos, peatones y semáforos, sino que también podemos inferir qué autos están circulando, cuáles están estacionados junto a la acera y hacia dónde se dirigen los peatones. Esta capacidad de inferir información dinámica a partir de una sola imagen es precisamente lo que la visión por computadora busca lograr.

\begin{figure}[H]
\centering
\includegraphics[width=0.85\textwidth]{slides-images/slide-03.jpg}
\caption{El objetivo de la visión por computadora: descubrir a partir de una imagen qué objetos hay en la escena, dónde están, qué acción está ocurriendo, y predecir e inferir eventos futuros. \protect\footnotemark}
\end{figure}
\footnotetext{Diapositiva 3.}

\begin{knowledgebox}{¿Por qué es tan difícil la visión por computadora?}
El sistema visual humano, tras millones de años de evolución, es capaz de realizar sin esfuerzo tareas como el reconocimiento de objetos y la comprensión de escenas. Pero para una computadora, incluso un mismo objeto puede presentar una apariencia radicalmente distinta a nivel de píxel debido a factores como la variación del punto de vista (viewpoint variation), la variación de escala (scale variation), la deformación (deformation), la oclusión (occlusion), las condiciones de iluminación (illumination conditions), el desorden del fondo (background clutter) y la variación intraclase (intra-class variation). Todos estos factores hacen de la visión por computadora un problema sumamente desafiante.
\end{knowledgebox}

\subsection{Aplicaciones e impacto de la visión por computadora}

El aprendizaje profundo está impulsando una revolución enorme en el campo de la visión por computadora, con aplicaciones que están en todas partes:

\begin{figure}[H]
\centering
\includegraphics[width=0.9\textwidth]{slides-images/slide-04.jpg}
\caption{Aplicaciones amplias de la visión por computadora: robótica, asistencia para la accesibilidad, biomedicina, conducción autónoma y computación móvil. \protect\footnotemark}
\end{figure}
\footnotetext{Diapositiva 4.}

\begin{itemize}
    \item \textbf{Detección y reconocimiento facial}: no solo detecta la presencia de un rostro, sino que también analiza puntos clave faciales y microexpresiones. Es una de las primeras tareas de visión por computadora que alcanzó un uso amplio.
    \item \textbf{Conducción autónoma}: el laboratorio de Amini en el MIT construyó en su momento un sistema de conducción autónoma de extremo a extremo capaz de conducir de manera autónoma en caminos nunca antes vistos utilizando solo entrada visual. A diferencia de la mayoría de los autos autónomos, que usan un pipeline modularizado (varios modelos más mapas predefinidos), el método de extremo a extremo usa un solo modelo que va directamente de la imagen al ángulo del volante.
    \item \textbf{Imágenes médicas}: las CNN ya superan la exactitud de radiólogos especializados en la detección de cáncer de mama, y pueden detectar lesiones que los médicos humanos pasan por alto. Además, se aplican ampliamente en la detección de cáncer de piel, el diagnóstico por imágenes de COVID-19 y otros campos.
    \item \textbf{Asistencia para la accesibilidad}: por ejemplo, el proyecto Google Project Guideline, que ayuda a personas con discapacidad visual a correr de manera independiente.
    \item \textbf{Computación móvil}: en el teléfono de cada persona se ejecutan una gran cantidad de modelos de visión por computadora.
\end{itemize}

\begin{figure}[H]
\centering
\begin{subfigure}[b]{0.48\textwidth}
\centering
\includegraphics[width=\textwidth]{slides-images/slide-05.jpg}
\caption{Detección facial y reconocimiento de puntos clave}
\end{subfigure}
\hfill
\begin{subfigure}[b]{0.48\textwidth}
\centering
\includegraphics[width=\textwidth]{slides-images/slide-07.jpg}
\caption{Visión por computadora en imágenes médicas}
\end{subfigure}
\caption{Aplicaciones representativas de la visión por computadora. \protect\footnotemark}
\end{figure}
\footnotetext{Diapositivas 5 y 7.}

\subsection{Resumen de la sección}

La visión por computadora consiste en dotar a las computadoras de la capacidad de comprender el mundo visual. Aunque para el ser humano resulta trivial, para la computadora es sumamente desafiante debido a la enorme variabilidad en la apariencia de los objetos. El auge del aprendizaje profundo está transformando este campo de manera radical, y permite que las computadoras alcancen e incluso superen el nivel humano en muchas tareas visuales.

\newpage
